\section{A weighted arctangent integral in classical trilogarithms}
\label{xid:wat:sec}

The canonical manuscript, in the paragraph following
\texttt{cleo:eq:lintanh}, considers the weighted analogue of the
Reshetnikov--Lee bilinear integral,
\begin{equation}\label{xid:wat:eq:K}
 K(p,q)=\int_0^{\pi/2}
 \frac{\arctan(p\sin x)\arctan(q\sin x)}{\sin x}\,dx,
 \qquad p,q\in\mathbb R.
\end{equation}
It explains why the unweighted reduction to one Legendre chi function
does not carry over, but gives no finite evaluation for this family.
We supply a two-term formula in a single even function built from
classical dilogarithms and trilogarithms at real arguments in $(0,1]$.
We also give a normalized ordinary antiderivative and an
arbitrary-order Euler integration ladder. The proof of the first
result uses parameter differentiation and hyperbolic addition;
the original unweighted proof of Lee~\cite{wat:Lee} also uses
hyperbolic parameters. The elementary logarithmic reduction used
later is part of standard weight-three polylogarithm calculus.
Its underlying two-logarithm kernel was treated for real parameters
by David H.~\cite{wat:DavidH}; the polarization method is not new.
The contribution relative to the inspected repository is the explicit
evaluation of its weighted two-parameter family, with real domains,
normalizations, and branch prescriptions proved. Historical priority
for this exact weighted evaluation has not been established.

The integral in \eqref{xid:wat:eq:K} is ordinary and convergent. Its integrand
is $pq\,x+O(x^3)$ at zero. The principal real arctangent is used
throughout. Principal complex logarithms and classical polylogarithms
are needed only for the more general ordinary primitive. Their series and integral
conventions agree with DLMF~\S25.12~\cite{wat:DLMF}.

\subsection{The hyperbolic addition formula}

Write
\[
 \chi_j(z)=\frac{\operatorname{Li}_j(z)-\operatorname{Li}_j(-z)}2
          =\sum_{k=0}^{\infty}\frac{z^{2k+1}}{(2k+1)^j}
 \qquad (|z|<1)
\]
for the Legendre chi function. The values used in the next theorem
are real and lie in its ordinary convergent series domain, including
the endpoints $\chi_2(1)=\pi^2/8$ and
$\chi_3(1)=7\zeta(3)/8$.

\begin{theorem}[A real two-term weighted evaluation]
\label{xid:wat:thm:main}
Define the even function
\begin{equation}\label{xid:wat:eq:Fclosed}
 \mathfrak F(s)=\frac{\pi^2|s|}{8}-\frac74\zeta(3)
    +|s|\chi_2(e^{-|s|})+2\chi_3(e^{-|s|}),
 \qquad s\in\mathbb R.
\end{equation}
Its value at zero is $\mathfrak F(0)=0$. The apparent lack of
smoothness from the absolute values is removable: $\mathfrak F$ is
real analytic on $\mathbb R$, and it is equivalently characterized by
\begin{equation}\label{xid:wat:eq:Fode}
 \mathfrak F''(s)=\frac{s}{2\sinh s},\qquad
 \mathfrak F''(0)=\frac12,\qquad
 \mathfrak F(0)=\mathfrak F'(0)=0.
\end{equation}
For every $p,q\in\mathbb R$, put
$\alpha=\operatorname{arcsinh}p$ and
$\beta=\operatorname{arcsinh}q$. Then
\begin{equation}\label{xid:wat:eq:main}
 \boxed{\displaystyle
 K(p,q)=\mathfrak F(\alpha+\beta)-\mathfrak F(\alpha-\beta).}
\end{equation}
In particular the formula includes both diagonals $p=\pm q$,
either zero parameter, and arbitrary signs without branch changes.
\end{theorem}

\begin{proof}
First, for $s>0$ use
$d\chi_j(e^{-s})/ds=-\chi_{j-1}(e^{-s})$ to obtain
\[
 \mathfrak F'(s)=\frac{\pi^2}{8}
              -\chi_2(e^{-s})-s\chi_1(e^{-s}),
 \qquad
 \mathfrak F''(s)=s\chi_0(e^{-s})
                =\frac{s}{2\sinh s}.
\]
The limits at zero follow from the endpoint values of $\chi_2$ and
$\chi_3$ and from
$s\chi_1(e^{-s})=\tfrac{s}{2}\log\coth(s/2)\to0$.
The unique twice-normalized primitive of the real analytic even
function $s/(2\sinh s)$ is real analytic and even. Hence
\eqref{xid:wat:eq:Fclosed} agrees with this primitive on the whole real
line; this proves all assertions in \eqref{xid:wat:eq:Fode}.

Differentiate \eqref{xid:wat:eq:K} once in each parameter. On compact
real parameter sets the differentiated integrand is bounded by an
integrable function, and therefore
\begin{align*}
 K_{pq}(p,q)
 &=\int_0^{\pi/2}
       \frac{\sin x\,dx}
       {(1+p^2\sin^2x)(1+q^2\sin^2x)}\\
 &=\frac{p\operatorname{arcsinh}p/\sqrt{1+p^2}
       -q\operatorname{arcsinh}q/\sqrt{1+q^2}}{p^2-q^2}
       \qquad(p^2\ne q^2).
\end{align*}
For the second equality use partial fractions and
\[
 \int_0^{\pi/2}\frac{\sin x}{1+p^2\sin^2x}\,dx
       =\frac{\operatorname{arcsinh}p}{p\sqrt{1+p^2}},
\]
whose zero-parameter value is $1$ and which follows by
$v=\cos x$.

Let $\widetilde K(\alpha,\beta)=K(\sinh\alpha,\sinh\beta)$.
Multiplication by $\cosh\alpha\cosh\beta$ and the addition formulas
for $\sinh$ now give
\begin{align}\label{xid:wat:eq:hyperbolicmixed}
 \widetilde K_{\alpha\beta}
 &=\frac{\alpha\sinh\alpha\cosh\beta
        -\beta\sinh\beta\cosh\alpha}
       {\sinh^2\alpha-\sinh^2\beta}\notag\\
 &=\frac12\left\{
      \frac{\alpha+\beta}{\sinh(\alpha+\beta)}
     +\frac{\alpha-\beta}{\sinh(\alpha-\beta)}\right\}.
\end{align}
Both quotients on the last line have their removable value $1$ at
zero. Thus the equality holds at $\alpha=\pm\beta$ by continuity.
By \eqref{xid:wat:eq:Fode}, this is the mixed derivative of
$\mathfrak F(\alpha+\beta)-\mathfrak F(\alpha-\beta)$.
Both this expression and $\widetilde K$ vanish on $\alpha=0$ and
$\beta=0$. Integration of their equal mixed derivatives over the
oriented rectangle from $(0,0)$ to $(\alpha,\beta)$ proves
\eqref{xid:wat:eq:main} for every real $\alpha,\beta$.
\end{proof}

For algebraic real $p,q$, the two arguments
$e^{-|\alpha+\beta|}$ and $e^{-|\alpha-\beta|}$ are algebraic.
Indeed $e^\alpha=p+\sqrt{1+p^2}$ and
$e^\beta=q+\sqrt{1+q^2}$ are positive algebraic numbers, and taking
the absolute value in either exponent selects an algebraic number
or its reciprocal. Thus \eqref{xid:wat:eq:main} expresses every such
integral by classical polylogarithms at real algebraic arguments and
logarithms of positive algebraic numbers. It asserts no independence
or minimality of the resulting constants.

The diagonal specializes particularly simply:
\begin{equation}\label{xid:wat:eq:Kdiagonal}
 K(p,p)=\mathfrak F(2\operatorname{arcsinh}p).
\end{equation}
For $\ell=\log(1+\sqrt2)$ and $r=3-2\sqrt2$, this gives the exact
identity
\begin{equation}\label{xid:wat:eq:K11}
 \boxed{\displaystyle
 \int_0^{\pi/2}\frac{\arctan^2(\sin x)}{\sin x}\,dx
 =\frac{\pi^2\ell}{4}-\frac74\zeta(3)
       +2\ell\chi_2(r)+2\chi_3(r).}
\end{equation}
Its numerical value is
$0.718064319741103407036328710693\ldots$; the decimal only
identifies the normalization, while the theorem proves the identity.

An exact weighted counterpart of the manuscript's Cleo integral
also follows. Set $\alpha=\log(1+\sqrt2)$ and
$\beta=\log\bigl((1+\sqrt5)/2\bigr)$. Then
\begin{align}\label{xid:wat:eq:weightedCleo}
 &\int_0^{\pi/2}\frac1{\sin x}
    \arctan^2\!\left(\frac{6\sin x}{3+\cos2x}\right)dx\notag\\
 &\qquad=\mathfrak F(2\alpha)
          +2\mathfrak F(\alpha+\beta)
          -2\mathfrak F(\alpha-\beta)
          +\mathfrak F(2\beta).
\end{align}
To verify the reduction put $s=\sin x\in[0,1]$. Since
$\arctan s+\arctan(s/2)<\pi/2$, the tangent addition formula gives
\[
 \arctan\frac{6\sin x}{3+\cos2x}
       =\arctan(\sin x)+\arctan(\tfrac12\sin x).
\]
Squaring and integrating yields
$K(1,1)+2K(1,1/2)+K(1/2,1/2)$, proving
\eqref{xid:wat:eq:weightedCleo} by \eqref{xid:wat:eq:main}.

\subsection{The weighted inverse-hyperbolic companion}

The same parameter method has a trigonometric form. It gives a
second exact family, including convergent improper endpoint values.

\begin{theorem}[An inverse-hyperbolic bilinear evaluation]
\label{xid:wat:thm:atanh}
For $|r|,|s|\leq1$ define
\[
 L(r,s)=\int_0^{\pi/2}
    \frac{\operatorname{arctanh}(r\sin x)
          \operatorname{arctanh}(s\sin x)}{\sin x}\,dx,
\]
using real inverse hyperbolic tangents and interpreting an endpoint
with $|r|=1$ or $|s|=1$ as an improper integral. Put
\begin{equation}\label{xid:wat:eq:Gclosed}
 \mathfrak G(t)=\frac74\zeta(3)
     -t\operatorname{Im}\chi_2(e^{it})
     -2\operatorname{Re}\chi_3(e^{it}),\qquad |t|\leq\pi.
\end{equation}
The chi values on the unit circle are defined by their absolutely
convergent series. Then $\mathfrak G$ is even, continuous on
$[-\pi,\pi]$, and real analytic on $(-\pi,\pi)$, with
\begin{equation}\label{xid:wat:eq:Gode}
 \mathfrak G''(t)=\frac{t}{2\sin t},\qquad
 \mathfrak G''(0)=\frac12,\qquad
 \mathfrak G(0)=\mathfrak G'(0)=0.
\end{equation}
For $a=\arcsin r$, $b=\arcsin s$, including either sign of either
parameter,
\begin{equation}\label{xid:wat:eq:atanhmain}
 \boxed{\displaystyle
 L(r,s)=\mathfrak G(a+b)-\mathfrak G(a-b).}
\end{equation}
In particular,
\begin{equation}\label{xid:wat:eq:atanh11}
 \boxed{\displaystyle
 \int_0^{\pi/2}\frac{\operatorname{arctanh}^2(\sin x)}{\sin x}\,dx
       =\frac72\zeta(3).}
\end{equation}
\end{theorem}

\begin{proof}
Conjugation of the chi series proves evenness and endpoint
continuity of \eqref{xid:wat:eq:Gclosed}. For $0<t<\pi$,
polylogarithm differentiation gives
\[
 \mathfrak G'(t)=\operatorname{Im}\chi_2(e^{it})
                    -t\operatorname{Re}\chi_1(e^{it}),
 \qquad
 \mathfrak G''(t)=t\operatorname{Im}\chi_0(e^{it})
                 =\frac{t}{2\sin t}.
\]
Here $\operatorname{Re}\chi_1(e^{it})=\tfrac12\log\cot(t/2)$,
so the limit of the first derivative at zero is zero. Also
$\mathfrak G(0)=0$ by $\chi_3(1)=7\zeta(3)/8$.
The normalized even primitive of $t/(2\sin t)$ is real analytic
on $(-\pi,\pi)$, proving \eqref{xid:wat:eq:Gode}.

For $|r|,|s|<1$, differentiation and partial fractions give
\[
 L_{rs}(r,s)=
 \frac{r\arcsin r/\sqrt{1-r^2}
       -s\arcsin s/\sqrt{1-s^2}}{r^2-s^2}
 \qquad(r^2\ne s^2).
\]
The underlying elementary integral is
$\int_0^{\pi/2}\sin x/(1-r^2\sin^2x)\,dx
 =\arcsin r/(r\sqrt{1-r^2})$, again obtained by $v=\cos x$.
For $\widetilde L(a,b)=L(\sin a,\sin b)$ the sine addition formulas
therefore yield
\[
 \widetilde L_{ab}(a,b)
   =\frac12\left\{
       \frac{a+b}{\sin(a+b)}+\frac{a-b}{\sin(a-b)}\right\},
 \qquad |a|,|b|<\pi/2.
\]
The zero quotients take their removable value $1$. This is the
mixed derivative of the right side of \eqref{xid:wat:eq:atanhmain},
and both sides vanish on the two axes. Rectangle integration proves
the identity in the open parameter square.

For all $|r|,|s|\leq1$ the absolute value of the integrand is bounded
by $\operatorname{arctanh}^2(\sin x)/\sin x$. This majorant is
$O(x)$ at zero and $O(\log^2(\pi/2-x))$ at the other endpoint,
so it is integrable. Dominated convergence extends the equality
to the closed square. Finally
$\chi_3(-1)=-7\zeta(3)/8$ and
$\operatorname{Im}\chi_2(-1)=0$ give
$\mathfrak G(\pi)=7\zeta(3)/2$, proving
\eqref{xid:wat:eq:atanh11}.
The endpoint value also follows directly from
$u=\operatorname{arctanh}(\sin x)$:
\[
 \int_0^\infty\frac{u^2}{\sinh u}\,du
    =4\sum_{k=0}^\infty\frac1{(2k+1)^3}
    =\frac72\zeta(3).
\]
\end{proof}

For algebraic $r,s\in[-1,1]$, the chi arguments in this theorem are
algebraic points on the unit circle, since
$e^{ia}=\sqrt{1-r^2}+ir$ and
$e^{ib}=\sqrt{1-s^2}+is$. The theorem has been proved directly over
real parameters; it does not depend on an unstated continuation of
the arctangent formula.

\subsection{The classical logarithmic kernel}

For $A,B>0$, we will need only the two types of kernel
\begin{equation}\label{xid:wat:eq:Ldef}
 \mathcal L(a,b)=\int_0^1
        \frac{\log(1-at)\log(1-bt)}{t}\,dt,
 \qquad (a,b)=(iA,iB)\text{ or }(iA,-iB).
\end{equation}
These integrals have no singularity on their integration path. Define
\begin{equation}\label{xid:wat:eq:P}
 P_c(y)=\log^2y\log(1-y/c)
       +2\log y\operatorname{Li}_2(y/c)
       -2\operatorname{Li}_3(y/c).
\end{equation}
The limiting convention is $P_1(1)=-2\zeta(3)$. In each other use of
$P_c(1)$ below, $1/c\notin[1,\infty)$, and thus
$P_c(1)=-2\operatorname{Li}_3(1/c)$ directly.

\begin{lemma}[A finite principal-branch logarithmic reduction]
\label{xid:wat:lem:L}
If $a=iA$, $b=iB$, $A>B>0$, or if $a=iA$, $b=-iB$, $A,B>0$, set
$r=(1-a)/(1-b)$. Then
\begin{equation}\label{xid:wat:eq:Lclosed}
 \boxed{\displaystyle
 \mathcal L(a,b)=\zeta(3)+\operatorname{Li}_3(b/a)
 +\frac{P_1(1-a)+P_1(1-b)-P_1(r)+P_{a/b}(r)}2.}
\end{equation}
For the same-sign case with $A<B$, interchange $a$ and $b$ before
applying this expression. At the diagonal,
\begin{equation}\label{xid:wat:eq:Ldiagonal}
 \mathcal L(a,a)=P_1(1-a)+2\zeta(3).
\end{equation}
All quantities in these prescriptions use principal branches; no
unstated continuation across a positive-real polylogarithm cut is needed.
\end{lemma}

\begin{proof}
Differentiation using
$d\operatorname{Li}_j(z)/dz=\operatorname{Li}_{j-1}(z)/z$ and
$\operatorname{Li}_1(z)=-\log(1-z)$ gives
\begin{equation}\label{xid:wat:eq:Pderivative}
 \frac{d}{dy}P_c(y)=\frac{\log^2y}{y-c}.
\end{equation}
The four extraneous terms cancel in pairs. Consequently
\[
 \int_0^1\frac{\log^2(1-at)}t\,dt
 =P_1(1-a)-P_1(1).
\]
For the cross term use the polarization identity
\[
 2\log(1-at)\log(1-bt)
 =\log^2(1-at)+\log^2(1-bt)
  -\log^2\frac{1-at}{1-bt}.
\]
Here the logarithm of the quotient equals the difference of the two
principal logarithms on the stated paths. Indeed, in the same-sign
case its argument is $-\arctan(At)+\arctan(Bt)\in(-\pi/2,0)$;
in the opposite-sign case it is
$-\arctan(At)-\arctan(Bt)\in(-\pi,0)$.

With $y=(1-at)/(1-bt)$ and $c=a/b$, elementary differentiation yields
\[
 \frac{dt}{t}=\left(\frac1{y-1}-\frac1{y-c}\right)dy.
\]
The quotient-log square therefore has integral
\[
 P_1(r)-P_1(1)-P_c(r)+P_c(1).
\]
Substitution into polarization proves \eqref{xid:wat:eq:Lclosed} and also
the diagonal formula.

For completeness, the branch check for the additional polylogarithms
is explicit. In the same-sign case $A>B$,
\[
 y(t)=\frac{1+ABt^2+i(B-A)t}{1+B^2t^2},
 \qquad
 \operatorname{Re}\frac{y(t)}c
 =\frac{B/A+B^2t^2}{1+B^2t^2}<1.
\]
Thus $y(t)$ is in the lower half-plane for $t>0$, while $y(t)/c$
starts in $(0,1)$ and avoids $[1,\infty)$.
In the opposite-sign case $y(t)$ is in the lower half-plane and
$y(t)/c$, with $c<0$, is in the upper half-plane, starting at
$-B/A<0$. These observations justify every use of
\eqref{xid:wat:eq:Pderivative} with the stated principal values. The
endpoint $y=1$ is removable in $P_1$ because
$\log^2y\log(1-y)\to0$ and
$\log y\operatorname{Li}_2(y)\to0$.
\end{proof}

Define the real symmetric kernel
\begin{equation}\label{xid:wat:eq:J}
 \mathcal J(A,B)=\frac12\operatorname{Re}
       \{\mathcal L(iA,-iB)-\mathcal L(iA,iB)\},
 \qquad A,B>0.
\end{equation}
Lemma~\ref{xid:wat:lem:L} makes this a finite expression in classical
trilogarithms, dilogarithms, logarithms, and $\zeta(3)$.
Conjugation of the logarithms gives the simpler integral meaning
\begin{equation}\label{xid:wat:eq:Jintegral}
 \mathcal J(A,B)=\int_0^1
              \frac{\arctan(At)\arctan(Bt)}t\,dt.
\end{equation}
To see the factor $1/2$, write
$\log(1-iAt)=\frac12\log(1+A^2t^2)-i\arctan(At)$;
the real part of the difference in \eqref{xid:wat:eq:J} is exactly twice
the product of arctangents.

\subsection{The ordinary primitive at a general endpoint}

\begin{theorem}[Finite classical evaluation and a normalized primitive]
\label{xid:wat:thm:primitive}
For $p,q\in\mathbb R$ put
\[
 a=e^{\operatorname{arcsinh}p}=p+\sqrt{1+p^2}>0,
 \qquad
 b=e^{\operatorname{arcsinh}q}=q+\sqrt{1+q^2}>0.
\]
For $0<\theta<\pi$ let $T=\tan(\theta/2)$. Then
\begin{align}\label{xid:wat:eq:primitive}
 &\int_0^\theta
   \frac{\arctan(p\sin x)\arctan(q\sin x)}{\sin x}\,dx
       \notag\\
 &\quad=\mathcal J(aT,bT)-\mathcal J(aT,T/b)
           -\mathcal J(T/a,bT)+\mathcal J(T/a,T/b).
\end{align}
In particular,
\begin{equation}\label{xid:wat:eq:fourkernel}
 \boxed{K(p,q)=\mathcal J(a,b)-\mathcal J(a,b^{-1})
       -\mathcal J(a^{-1},b)+\mathcal J(a^{-1},b^{-1}).}
\end{equation}
Together with Lemma~\ref{xid:wat:lem:L}, these are finite formulas using
only classical polylogarithms of orders at most three. They include
negative and zero values of $p,q$ and the diagonals $p=\pm q$.
\end{theorem}

\begin{proof}
The half-angle substitution $t=\tan(x/2)$ gives
\[
 \sin x=\frac{2t}{1+t^2},\qquad \frac{dx}{\sin x}=\frac{dt}{t}.
\]
Since $a-a^{-1}=2p$,
\begin{equation}\label{xid:wat:eq:atandifference}
 \arctan\frac{2pt}{1+t^2}
       =\arctan(at)-\arctan(t/a),\qquad t>0.
\end{equation}
The tangents of the two sides coincide. Both sides lie strictly
between $-\pi/2$ and $\pi/2$ because the two arctangents on the right
belong to $(0,\pi/2)$. Hence the equality has no additive multiple
of $\pi$, including when $p<0$.
Apply \eqref{xid:wat:eq:atandifference} for $p$ and $q$, expand the
product, and rescale each integral over $[0,T]$ to $[0,1]$.
Equation~\eqref{xid:wat:eq:Jintegral} gives \eqref{xid:wat:eq:primitive}.
The choice $\theta=\pi/2$ gives $T=1$ and proves
\eqref{xid:wat:eq:fourkernel}.
\end{proof}

The right side of \eqref{xid:wat:eq:primitive}, plus an arbitrary constant,
is an ordinary antiderivative in $\theta$ on $(0,\pi)$. Its constant
has been fixed by the limit zero at $\theta=0$. For algebraic $p,q$
and algebraic $T>0$, every argument in the finite formula is algebraic.
This includes rational multiples $\theta/\pi$ with $0<\theta<\pi$.
The statement is an explicit upper bound on the kinds of functions
needed for the expression; it asserts no independence or minimality
of its special values.

\subsection{An elementary differential certificate}

The elementary computation in the proof of
Theorem~\ref{xid:wat:thm:main} supplies a differential certificate for
both finite formulas, including their diagonal limits.

\begin{proposition}[Mixed derivative]
\label{xid:wat:prop:derivative}
For real $p,q$ with $p^2\ne q^2$,
\begin{equation}\label{xid:wat:eq:mixedderivative}
 \partial_p\partial_qK(p,q)
 =\frac{p\operatorname{arcsinh}p/\sqrt{1+p^2}
       -q\operatorname{arcsinh}q/\sqrt{1+q^2}}{p^2-q^2}.
\end{equation}
At $p^2=q^2$ the continuous value, for $p\ne0$, is
\begin{equation}\label{xid:wat:eq:mixeddiagonal}
 \frac1{2(1+p^2)}
 +\frac{\operatorname{arcsinh}p}{2p(1+p^2)^{3/2}},
\end{equation}
and the value at $(0,0)$ is $1$.
\end{proposition}

\begin{proof}
Equation~\eqref{xid:wat:eq:mixedderivative} was proved above. Take its
divided-difference limit by differentiating the numerator as a
function of $p^2$ to obtain \eqref{xid:wat:eq:mixeddiagonal}.
Its continuous limit at zero is $1$.
\end{proof}

Since $K(p,0)=K(0,q)=0$, its mixed derivative determines $K$ uniquely
by integration over the rectangle with corners $(0,0)$ and $(p,q)$.
This supplies an elementary differential certificate for any further
compressed trilogarithm representation of the same family.

\subsection{Every Euler integration order}

The half-angle reduction is not confined to weight three. For
integers $m\geq1$ set
\begin{equation}\label{xid:wat:eq:Lm}
 \mathcal L_m(a,b)=\frac1{(m-1)!}\int_0^1
    \frac{(-\log t)^{m-1}\log(1-at)\log(1-bt)}t\,dt.
\end{equation}
For the imaginary arguments used below this is an ordinary convergent
integral. The convention for double polylogarithms is
\[
 \operatorname{Li}_{r,1}(x,y)
 =\sum_{n>k\geq1}\frac{x^ny^k}{n^r k}.
\]

\begin{theorem}[A depth-two Euler primitive ladder]
\label{xid:wat:thm:ladder}
For $m\geq1$,
\begin{equation}\label{xid:wat:eq:LmMPL}
 \mathcal L_m(a,b)
 =\operatorname{Li}_{m+1,1}(a,b/a)
  +\operatorname{Li}_{m+1,1}(b,a/b).
\end{equation}
This is initially a series identity for $|a|,|b|<1$ and nonzero $a,b$.
For purely imaginary nonzero $a,b$, continue the two arguments
simultaneously along $(a,b)\mapsto(\lambda a,\lambda b)$,
$0<\lambda\leq1$. This radial prescription gives the value in
\eqref{xid:wat:eq:Lm}; no singularity is crossed.

Define
\[
 \mathcal J_m(A,B)=\frac12\operatorname{Re}
     \{\mathcal L_m(iA,-iB)-\mathcal L_m(iA,iB)\},
 \qquad A,B>0,
\]
and abbreviate
\[
 D_p(t)=\arctan(at)-\arctan(t/a)
       =\arctan\frac{2pt}{1+t^2}.
\]
Set
\begin{align}\label{xid:wat:eq:Imclosed}
 \mathcal I_m(T;p,q)
 ={}&\mathcal J_m(aT,bT)-\mathcal J_m(aT,T/b)\notag\\
   &-\mathcal J_m(T/a,bT)+\mathcal J_m(T/a,T/b).
\end{align}
Then
\begin{equation}\label{xid:wat:eq:Imint}
 \mathcal I_m(T;p,q)
 =\frac1{(m-1)!}\int_0^T
    \log^{m-1}(T/t)D_p(t)D_q(t)\,\frac{dt}{t}.
\end{equation}
Writing $\Theta_T=T\,d/dT$,
\begin{equation}\label{xid:wat:eq:Eulerladder}
 \Theta_T\mathcal I_m=\mathcal I_{m-1}\quad(m\geq2),
 \qquad
 \Theta_T\mathcal I_1=D_p(T)D_q(T).
\end{equation}
Thus every fixed integration order has a finite expression in double
polylogarithms of weight $m+2$ and depth at most two.
\end{theorem}

\begin{proof}
Expand the logarithms for $|a|,|b|<1$ and integrate term by term.
Absolute convergence gives
\[
 \mathcal L_m(a,b)
 =\sum_{j,k\geq1}\frac{a^jb^k}{jk(j+k)^m}.
\]
Use $1/(jk)=(1/j+1/k)/(j+k)$ and write $n=j+k$.
The two resulting sums are exactly those in
\eqref{xid:wat:eq:LmMPL}. Along the indicated imaginary radial paths,
$1-\lambda at$ and $1-\lambda bt$ never vanish for $0\leq t\leq1$;
analytic continuation gives the asserted extension. The real-part
calculation establishing \eqref{xid:wat:eq:Jintegral} is unchanged by
the logarithmic weight. The four-term expansion and the scaling
$t=Ts$ prove \eqref{xid:wat:eq:Imint}. Differentiating its normalized
Cauchy integration kernel gives \eqref{xid:wat:eq:Eulerladder}.
Near $t=0$ the product $D_p(t)D_q(t)$ is $O(t^2)$, which
justifies every differentiation and fixes the primitive constants.
\end{proof}

One useful diagonal identity follows without another integration:
\begin{equation}\label{xid:wat:eq:diagonalMPL}
 \mathcal J_m(A,A)
 =2^{-m-1}\operatorname{Li}_{m+1,1}(-A^2,1)
   -2\operatorname{Re}\operatorname{Li}_{m+1,1}(iA,1).
\end{equation}
Indeed, polarization of
$\log(1-at)+\log(1+at)=\log(1-a^2t^2)$ gives
\[
 \mathcal L_m(a,-a)
 =2^{-m}\operatorname{Li}_{m+1,1}(a^2,1)
   -\operatorname{Li}_{m+1,1}(a,1)
   -\operatorname{Li}_{m+1,1}(-a,1),
\]
where $t\mapsto t^2$ accounts for $2^{-m}$. Combine this with
$\mathcal L_m(a,a)=2\operatorname{Li}_{m+1,1}(a,1)$.
The arguments in \eqref{xid:wat:eq:diagonalMPL} lie off the positive-real
cut, so the same radial continuation is unambiguous.

\subsection{An exact obstruction to the one-product ansatz}

The earlier paragraph's failure of the unweighted factorization can
be made into a very short exact statement. Set
\[
 u(p)=\frac{p}{1+\sqrt{1+p^2}},\qquad u(0)=0.
\]
The unweighted integral is a function of $u(p)u(q)$. The weighted
integral has no such analytic factorization, even into an unspecified
one-variable function.

\begin{proposition}\label{xid:wat:prop:nonfactor}
There is no function $G$ analytic near zero for which
$K(p,q)=G(u(p)u(q))$ holds for all sufficiently small real $p,q$.
\end{proposition}

\begin{proof}
For $|p|,|q|<1$, integration of the absolutely convergent arctangent
series gives
\begin{equation}\label{xid:wat:eq:doubleTaylor}
 K(p,q)=\sum_{j,k\geq0}
 \frac{(-1)^{j+k}p^{2j+1}q^{2k+1}}{(2j+1)(2k+1)}
 \frac{4^{j+k}}{(2j+2k+1)\binom{2j+2k}{j+k}}.
\end{equation}
In particular,
\[
 K(p,q)=pq-\frac29(p^3q+pq^3)
       +\text{terms of total degree at least six},
\]
whereas
\[
 u(p)u(q)=\frac{pq}{4}-\frac{p^3q+pq^3}{16}
       +\text{terms of total degree at least six}.
\]
The coefficient of $pq$ would force $G'(0)=4$, and hence the
coefficient of $p^3q$ would have to be $-1/4$, contradicting $-2/9$.
Terms of order at least two in $G$ cannot contain $p^3q$.
\end{proof}

This is an obstruction only to the stated functional ansatz. It says
nothing about independence of particular constants. Theorem~\ref{xid:wat:thm:main}
shows explicitly that a short formula does exist after changing to
the hyperbolic sum and difference of the parameters.
Unlike a bounded search for one $\chi_3$, the coefficient mismatch
proves the exact assertion for every analytic $G$.

\paragraph{Reproduction and next questions.}
The accompanying verification program checks
\eqref{xid:wat:eq:Pderivative} and the nonfactorization coefficients
symbolically, compares both classical formulas with direct real
quadrature, verifies the hyperbolic mixed derivative and the
ordinary primitive at several endpoints,
tests the inverse-hyperbolic family at signed and nearly saturated
endpoints, and checks the higher integration orders against independent nested
series in a common convergence disk. These calculations audit
normalizations and branch choices; the proofs above establish the
identities. Natural next tasks are to identify additional endpoints
at which the four-kernel ordinary primitive compresses to real
one-variable kernels, to seek analogous addition formulas for higher
Euler integration orders, and to determine which specializations
of that higher-weight ladder admit classical polylogarithm reductions.
The Gaussian $S_6$ and current $S_8$
candidates remain unresolved by this argument.
